\documentclass[crop=false]{standalone}
\usepackage{15150toolbox}
\usepackage{comment}
\usepackage{needspace}

% Toggle: set \soltrue to render solutions inline (blue boxes).
\newif\ifsol
\soltrue   % comment this line for the student version

\ifsol
  \newenvironment{solution}
    {\par\color{solution_color}\begin{framed}\textbf{Solution.}\par}
    {\end{framed}}
\else
  \excludecomment{solution}
\fi

% Lecture-slide macros that aren't defined in 15150toolbox.
\newcommand{\term}[1]{\textbf{\textit{#1}}}
\newcommand{\thmBox}[2]{\par\noindent\begin{mdframed}[backgroundcolor=blue!8,hidealllines=true]\textbf{Theorem#1.} #2\end{mdframed}}
\newcommand{\bcBox}[1]{\par\noindent\textbf{Base case.} #1\par}
\newcommand{\ihBox}[2]{\par\noindent\textbf{Inductive hypothesis#1.} #2\par}
\newcommand{\isBox}[2]{\par\noindent\textbf{Inductive step#1.} #2\par}

\newcommand{\quizproblem}[3]{%
  \needspace{5\baselineskip}%
  \stepcounter{problem}%
  \ifsol
    \subsection*{Problem \arabic{problem}. #1 \hfill \timing{#2} \quad (#3 pts)}%
  \else
    \subsection*{Problem \arabic{problem}. #1 \hfill (#3 pts)}%
  \fi
}

\newcommand{\subpts}[1]{\textbf{(#1 pts)}}

\newcommand{\answerspace}[1]{\ifsol\else\vspace{#1}\par\fi}

\printout{Quiz 2}

\begin{document}

\begin{center}
  {\Large \textbf{15-150 --- Quiz 2}} \\
  \vspace{4pt}
  Topics: \textit{Induction and Recursion}; \textit{Structural Induction and Tail Recursion}; \\
  \textit{Trees}; \textit{Asymptotic Analysis} \\
  \vspace{4pt}
  \textit{60 minutes.}
\end{center}

\vspace{1em}

\noindent\textbf{Name:} \hrulefill \hspace{2em} \textbf{Andrew ID:} \hrulefill

\vspace{1em}

% =============================================================================
% PROBLEM 1 --- Write contains
% =============================================================================

\quizproblem{A Tree Function}{6}{8}

The \code{tree} datatype is defined as:

\begin{codeblock}
  datatype tree = Empty | Node of tree * int * tree
\end{codeblock}

Write a function \code{contains} satisfying the following specification:

\spec
  {contains}
  {int * tree -> bool}
  {true}
  {\code{contains (n, T)} evaluates to \code{true} if \code{n} appears as
   the value at some \code{Node} in \code{T}, and \code{false} otherwise.}

\answerspace{6cm}

\begin{solution}
\begin{codeblock}
  fun contains (n : int, Empty : tree) : bool = false
    | contains (n, Node (L, x, R)) =
        n = x orelse contains (n, L) orelse contains (n, R)
\end{codeblock}

\textbf{Rubric:}
\begin{itemize}
  \item 3 pts: \code{Empty} $\to$ \code{false}
  \item 5 pts: \code{Node} case combines equality check with both
    recursive calls via \code{orelse}. $-1$ per missing piece (eq check,
    \code{L} recur, \code{R} recur). \code{andalso}: $-2$.
\end{itemize}
\code{case} on the tuple, reordered disjuncts, etc.\ are fine.
\end{solution}

% =============================================================================
% PROBLEM 2 --- Datatype types
% =============================================================================

\vspace{1em}
\quizproblem{Datatype Constructors}{6}{8}

Consider the following \code{datatype} declaration:

\begin{codeblock}
  datatype coin = Heads
                | Tails
                | Stack of coin * coin
\end{codeblock}

For each expression below, write its type if it is well-typed, or write
``ill-typed'' with a short reason.

\vspace{0.5em}

\begin{center}
\begin{tabular}{|l|p{6.5cm}|}
  \hline
  \textbf{Expression} & \textbf{Type, or ill-typed (why)} \\
  \hline
  \code{Stack}                                    & \\[2em] \hline
  \code{Stack (Heads, Stack (Tails, Heads))}      & \\[2em] \hline
  \code{Stack (Heads, 5)}                         & \\[2em] \hline
  \code{SOME Heads}                               & \\[2em] \hline
  \code{[Heads, 0]}                               & \\[2em] \hline
\end{tabular}
\end{center}

\begin{solution}
\begin{center}
\begin{tabular}{|l|l|}
  \hline
  \code{Stack}                                    & \code{coin * coin -> coin} \\ \hline
  \code{Stack (Heads, Stack (Tails, Heads))}      & \code{coin} \\ \hline
  \code{Stack (Heads, 5)}                         & ill-typed (\code{5 : int}, not \code{coin}) \\ \hline
  \code{SOME Heads}                               & \code{coin option} \\ \hline
  \code{[Heads, 0]}                               & ill-typed (heterogeneous list) \\ \hline
\end{tabular}
\end{center}

\textbf{Rubric:}
\begin{itemize}
  \item 2 pts each: \code{Stack}, \code{Stack (Heads, 5)}, \code{[Heads, 0]}
  \item 1 pt each: \code{Stack (Heads, Stack (Tails, Heads))}, \code{SOME Heads}
\end{itemize}
Ill-typed without a reason: half credit.
The likely error on \code{Stack} alone is writing \code{coin}.
\end{solution}

% =============================================================================
% PROBLEM 3 --- Strengthening the theorem
% =============================================================================

\clearpage
\quizproblem{Strengthening the Theorem}{6}{12}

Recall the tail-recursive reverse function \code{trev} and its wrapper
\code{rev} from lecture:

\begin{codeblock}
  fun trev ([] : int list, acc : int list) : int list = acc
    | trev (x::xs, acc) = trev (xs, x::acc)

  fun rev (L : int list) : int list = trev (L, [])
\end{codeblock}

Also recall the naive reverse \code{rev'} from lecture, which we will
use as the reference implementation:

\begin{codeblock}
  fun rev' ([] : int list) : int list = []
    | rev' (x::xs) = rev' xs @ [x]
\end{codeblock}

A student tries to prove the following theorem by structural induction
on \code{L}:

\thmBox{}{\, For all values \code{L : int list},
\code{rev L} $\eeq$ \code{rev' L}.}

The base case goes through fine. In the inductive case \code{L = x::xs},
the student assumes the IH \code{rev xs} $\eeq$ \code{rev' xs}, and
writes:

\begin{align*}
  \code{rev (x::xs)} \eeq\ &\code{trev (x::xs, [])}      \tag{def.\ of \code{rev}} \\
                    \eeq\ &\code{trev (xs, [x])}        \tag{clause 2 of \code{trev}} \\
                    \eeq\ &\ ???
\end{align*}

\begin{quote}
\emph{``...Now what do I do?''}
\end{quote}

\vspace{0.3em}

\textbf{(a)} \subpts{4} In \textbf{one sentence}, explain why the IH
above is too weak to finish this proof.

\answerspace{3cm}

\begin{solution}
The IH is about \code{trev (xs, [])}, but the expression we have is
\code{trev (xs, [x])} --- a different accumulator value.

\textbf{Rubric:} 4 pts for naming the accumulator-value mismatch.
``IH is wrong'' without saying why: 1 of 4.
\end{solution}

\vspace{0.3em}

\textbf{(b)} \subpts{8} Write a \textbf{stronger theorem} which (i)
implies the original theorem as a special case, and (ii) would give an
inductive hypothesis general enough to finish the proof. You do
\emph{not} need to write the proof, just the strengthened theorem
statement.

\answerspace{4cm}

\begin{solution}
For all values \code{L : int list} and all values \code{acc : int list},
\code{trev (L, acc)} $\eeq$ \code{rev' L @ acc}.
Original theorem is the \code{acc = []} case (since \code{rev' L @ []}
$\eeq$ \code{rev' L}).

\textbf{Rubric:}
\begin{itemize}
  \item 4 pts: quantifies over all \code{acc}
  \item 4 pts: RHS is \code{rev' L @ acc} (the accumulator gets
    \emph{appended}). \code{rev' L} alone: 0 on this bullet.
    \code{acc @ rev' L}: $-2$ (wrong order). \code{rev' L :: acc} or
    \code{rev' L + acc}: 0 --- doesn't typecheck.
\end{itemize}
Right quantifier, wrong RHS: 4 of 8.
\end{solution}

% =============================================================================
% PROBLEM 4 --- Structural induction: invert (invert T) eeq T
% =============================================================================

\clearpage
\quizproblem{An Induction Proof}{13}{12}

Consider the following function, which mirrors a tree by swapping every
left and right subtree:

\begin{codeblock}
  fun invert (Empty : tree) : tree = Empty
    | invert (Node (L, x, R)) = Node (invert R, x, invert L)
\end{codeblock}

\thmBox{}{\, For all values \code{T : tree},
\code{invert (invert T)} $\eeq$ \code{T}.}

You may freely use, without proof:
\begin{itemize}
  \item \textbf{Lemma 1.} \code{invert} is total.
\end{itemize}

\vspace{0.3em}

\textbf{Prove the theorem.} Use structural induction on \code{T}. Be sure
to state your inductive hypotheses explicitly and justify each step of
your calculation.

\answerspace{14cm}

\begin{solution}
Structural induction on \code{T : tree}.

\vspace{3pt}

\bcBox{\, \code{T = Empty}.}
\begin{align*}
  \code{invert (invert Empty)}
    \eeq\ &\code{invert Empty}   \tag{Clause 1 of \code{invert}}\\
    \eeq\ &\code{Empty}          \tag{Clause 1 of \code{invert}}
\end{align*}

\ihBox{}{\, For some \code{L, R : tree}, \code{invert (invert L)} $\eeq$
\code{L} and \code{invert (invert R)} $\eeq$ \code{R}.}

\isBox{}{\, \code{T = Node (L, x, R)}. WTS
\code{invert (invert (Node (L, x, R)))} $\eeq$ \code{Node (L, x, R)}.}

\begin{align*}
  &\code{invert (invert (Node (L, x, R)))} \\
  \eeq\ &\code{invert (Node (invert R, x, invert L))}
      \tag{Clause 2 of \code{invert}, \code{invert} total}\\
  \eeq\ &\code{Node (invert (invert L), x, invert (invert R))}
      \tag{Clause 2 of \code{invert}, \code{invert} total}\\
  \eeq\ &\code{Node (L, x, R)}
      \tag{IH twice}
\end{align*}

\textbf{Rubric:}
\begin{itemize}
  \item 2 pts: structural induction on \code{T} (not simple/strong)
  \item 2 pts: base case
  \item 2 pts: two IHs, quantified ``for some''. ``for all'': $-1$.
    Only one IH: $-1$.
  \item 4 pts: inductive step calculation. Missing totality citation:
    $-1$. Using IH only once: $-2$. Each justified equational step:
    1 pt.
  \item 2 pts: IHs applied correctly to land on \code{Node (L, x, R)}
\end{itemize}
\end{solution}

% =============================================================================
% PROBLEM 5 --- Work and span recurrences
% =============================================================================

\vspace{1em}
\quizproblem{Work and Span Recurrences}{7}{8}

Consider the following function on trees:

\begin{codeblock}
  fun addOne (Empty : tree) : tree = Empty
    | addOne (Node (L, x, R)) = Node (addOne L, x + 1, addOne R)
\end{codeblock}

Let $n$ denote the number of nodes in the input tree. For both parts
below, you may use $c_0, c_1, \ldots$ for unspecified constants.
\textbf{You do not need to solve the recurrence --- just write it down.}

\vspace{0.5em}

\textbf{(a)} \subpts{4} Write the recurrence for the \textbf{work} of
\code{addOne}, in terms of $n$, \textbf{assuming the tree is maximally
unbalanced} (i.e., the worst case).

\answerspace{3.5cm}

\begin{solution}
\begin{align*}
  W_{\code{addOne}}(0) &= c_0 \\
  W_{\code{addOne}}(n) &= W_{\code{addOne}}(n - 1) + W_{\code{addOne}}(0) + c_1
\end{align*}
Worst case: one subtree has $n-1$ nodes, the other has $0$. Simplified
form $W(n-1) + c$ also fine.

\textbf{Rubric:} 1 pt base, 3 pts recursive case. Balanced split
$W(n/2) + W(n/2) + c$: $-2$ (problem said worst case).
\end{solution}

\vspace{0.5em}

\textbf{(b)} \subpts{4} Now write the recurrence for the \textbf{span} of
\code{addOne}, \textbf{still assuming the tree is maximally unbalanced}.
Remember that, for span computations, tuples are evaluated in parallel.
Briefly comment on what this says about parallelism in this case.

\answerspace{4cm}

\begin{solution}
\begin{align*}
  S_{\code{addOne}}(0) &= c_0 \\
  S_{\code{addOne}}(n) &= \max\!\left(S_{\code{addOne}}(n-1),\,S_{\code{addOne}}(0)\right) + c_1 \\
                       &= S_{\code{addOne}}(n-1) + c_1
\end{align*}
The two recursive calls run in parallel (tuple components), so we take
$\max$ --- but on a maximally unbalanced tree, one subtree has $n-1$
nodes and the other has $0$, so $\max$ is just the big one. The
recurrence collapses to the same shape as the work recurrence:
\textbf{parallelism doesn't help on a maximally unbalanced tree.}

\textbf{Rubric:} 1 pt base, 2 pts recursive case (the
$\max(S(n-1), S(0))$ form, or simplified $S(n-1) + c$).
1 pt comment that span $=$ work here / parallelism doesn't help.
$S(n/2) + c$ (ignored that tree is unbalanced): $-2$.
$S(n-1) + S(0) + c$ (sum, not max): $-1$ --- gets the unbalanced split
right but missed the parallelism step.
\end{solution}

\end{document}
